% Generated draft: requires mathematical and source review.
\documentclass[11pt]{article}
\usepackage[margin=25mm]{geometry}
\usepackage{fontspec}
\setmainfont{Latin Modern Roman}
\usepackage{amsmath,amssymb,mathtools,unicode-math}
\setmathfont{Latin Modern Math}
\usepackage{enumitem,needspace,xurl,hyperref}
\setlength{\parindent}{0pt}
\setlength{\parskip}{5pt}
\setlength{\emergencystretch}{3em}
\newcommand{\N}{\mathbb{N}}
\newcommand{\Z}{\mathbb{Z}}
\newcommand{\Q}{\mathbb{Q}}
\newcommand{\R}{\mathbb{R}}
\newcommand{\C}{\mathbb{C}}
\newcommand{\F}{\mathbb{F}}
\newcommand{\A}{\mathbb{A}}
\newcommand{\PP}{\mathbb{P}}
\newcommand{\E}{\mathbb{E}}
\newcommand{\eps}{\varepsilon}
\newcommand{\dd}{\,\mathrm{d}}
\begin{document}
\section*{UnsolvedMath ID: 800796}
\textbf{Source problem code:} catalogue entry
\subsection{Definitions and mathematical statement}
Let \(T\) be a topological torus, that is, a surface homeomorphic to \(S^1\times S^1\). Interpret a simple loopless triangulation of \(T\) as a finite connected graph \(G\), with no loops or parallel edges, embedded cellularly in \(T\): every complementary face is an open disk, and the boundary of each face is a cycle of length three. This uses the conventional finite cellular interpretation of triangulation; these conventions are not explicitly stated in the supplied excerpt.

The embedded dual graph \(G^*\) has one vertex for each face of \(G\) and one edge \(e^*\) for each edge \(e\) of \(G\), joining the vertices corresponding to the faces on the two sides of \(e\). Distinct dual edges are retained even when they have the same endpoints. A proper \(3\)-edge-coloring assigns one of three colors to each edge so that the three edges incident with each dual vertex receive pairwise distinct colors.

\textbf{Conjecture 3 (supplied attribution: Kahle, 2019).} For every simple loopless triangulation \(G\) of \(T\), the dual graph \(G^*\) admits a proper \(3\)-edge-coloring. Equivalently, there exists a map
\[
 c:E(G)\longrightarrow\{1,2,3\}
\]
such that, for every triangular face \(f\) whose boundary edges are \(e_1,e_2,e_3\),
\[
 \{c(e_1),c(e_2),c(e_3)\}=\{1,2,3\}.
\]
The source introduces this conjecture with a statement that the preceding question remains open for genus \(g\) with \(1\leq g\leq4\), and describes the torus assertion as apparently open. The preceding general question is not supplied, so no conjecture for the other genera is formulated here. These are supplied status qualifications, not independently verified mathematical or literature assessments.
\subsection{Short English statement}
The conjecture asserts that every finite simple graph triangulating a torus has a dual whose edges can be colored with three colors, with all three colors appearing at each dual vertex. Equivalently, color the original edges so that every triangular face has three differently colored sides. The source describes this assertion as apparently open.
\subsection{Sources}
\begin{itemize}
\item Local UnsolvedMath registry, supplied complete record with permanent ID \(800796\), exact source problem code \texttt{catalogue entry}, and upstream revision \texttt{d371ee917ac9647fbb06cd6fc529deec6f3dda82}. The supplied registry status is \texttt{open}.
\item Supplied source attribution: Kahle, \emph{Branko Grunbaum in Many Dimensions: Open Problems} (2019), Conjecture 3. Supplied source URL: \url{https://arxiv.org/abs/1901.08622}. The record identifies the extraction method as \texttt{arxiv\_latex} and reports access on \texttt{2026-10-16}; this access information is supplied provenance and has not been independently verified.
\end{itemize}
\end{document}
